%【紙面設定】
%b4,横書き,2段
\documentclass[paper=b4j,landscape,twocolumn,fleqn]{jlreq}
%横書き用の調整
\special{papersize=\the\paperwidth,\the\paperheight}
%間に線をいれる
\setlength{\columnseprule}{0.4pt}
%ページ番号消す
\pagestyle{empty}
%ページ余白の調整
\usepackage[margin=10truemm]{geometry}
%画像挿入用のパッケージ
\usepackage[dvipdfmx]{graphicx}
%数学用のパッケージ
\usepackage{amsmath}
% 見出しの設定：行取りを1行にし、前後のスペースを極限まで削る
\ModifyHeading{section}{
  lines=1,
  before_space=10pt, % セクションの上の余白（適度に残すと見やすい）
  after_space=0pt    % セクションのしたの余白をゼロにする
}


%【本文】
\begin{document}
%1段落目
\section{五則演算(÷は/で表す.×は*で表す.)}
\begin{align*}
    5+8 &= & 9-4 &= & 2*3 &= & 6/2 &= & 5\%2 &= \\
    -9+4 &= & 3-10 &= & -4*5 &= & -10/5 &= & 10\%3 &= \\
    15+27 &= & 50-18 &= & 12*4 &= & 48/4 &= & 25\%4 &= \\
    -50+12 &= & 20-85 &= & 15*-3 &= & 100/-25 &= & 100\%7 &= \\
    1.2+3.8 &= & 4.5-1.2 &= & 0.5*6 &= & 4.8/2 &= & 50\%6 &= \\
    -4.5+1.2 &= & 0.8-3.5 &= & -1.2*4 &= & -7.2/0.9 &= & 80\%9 &= \\
    125+380 &= & 500-125 &= & 25*12 &= & 144/12 &= & 123\%10 &= \\
    -600+250 &= & 100-456 &= & 14*-10 &= & -117/9 &= & 250\%15 &= \\
    0.45+0.55 &= & 1.25-0.75 &= & 0.8*0.5 &= & 0.75/0.25 &= & 1000\%3 &= \\
    -12.5+7.5 &= & 1000-1500 &= & -15*1.2 &= & 102/-3 &= & 1024\%10 &= 
\end{align*}
\section{ルートと乗数の計算}
\begin{align*}
    2^2 &= & 2^2 * 2^1 &= & \sqrt{4} &= & \sqrt{4} * \sqrt{9} &= \\
    3^2 &= & 3^3 / 3^1 &= & \sqrt{9} &= & \sqrt{2} * \sqrt{8} &= \\
    2^4 &= & 2^2 * 2^2 &= & \sqrt{16} &= & \sqrt{12} / \sqrt{3} &= \\
    5^2 &= & 10^2 / 10^1 &= & \sqrt{25} &= & \sqrt{3} * \sqrt{12} &= \\
    4^3 &= & 2^3 * 2^2 &= & \sqrt{49} &= & \sqrt{18} / \sqrt{2} &= \\
    6^2 &= & 5^3 / 5^2 &= & \sqrt{64} &= & \sqrt{2} * \sqrt{32} &= \\
    10^2 &= & 2^2 * 2^4 &= & \sqrt{100} &= & \sqrt{50} / \sqrt{2} &= \\
    12^2 &= & 3^4 / 3^2 &= & \sqrt{144} &= & \sqrt{72} / \sqrt{2} &= 
\end{align*}
\section{分数の計算}
\begin{align*}
    \frac{1}{4} + \frac{1}{4} &= & \frac{4}{5} - \frac{1}{5} &= & \frac{1}{3} \times \frac{1}{2} &= & \frac{2}{5} \div 2 &= \\
    -\frac{2}{3} + \frac{1}{3} &= & \frac{1}{4} - \frac{3}{4} &= & -\frac{1}{2} \times \frac{1}{4} &= & \frac{1}{5} \div (-2) &= \\
    \frac{1}{2} + \frac{1}{3} &= & \frac{1}{2} - \frac{1}{4} &= & \frac{2}{3} \times \frac{3}{5} &= & \frac{3}{4} \div 2 &= \\
    -\frac{1}{2} + \frac{1}{6} &= & \frac{1}{6} - \frac{1}{2} &= & \frac{2}{5} \times \left(-\frac{1}{4}\right) &= & -\frac{3}{4} \div 3 &= \\
    \frac{3}{4} + \frac{1}{6} &= & \frac{5}{6} - \frac{1}{4} &= & \frac{5}{8} \times \frac{4}{15} &= & \frac{2}{3} \div \frac{1}{2} &= \\
    -\frac{1}{2} - \frac{1}{3} - \frac{1}{6} &= & \frac{1}{10} - \frac{4}{15} &= & -\frac{15}{16} \times \frac{4}{5} &= & 10 \div \left(-\frac{5}{2}\right) &= 
\end{align*}

%2段落目
\newpage
\section{方程式の計算(xを求めよ)}
\begin{align*}
    x + 5 &= 8 & 3x &= 12 & x + 10 &= 4 \\
    2x - 3 &= 7 & 15 - x &= 20 & 4x + 16 &= 0 \\
    3x + 4 &= x + 10 & 2x - 5 &= 5x + 7 & 4x - 8 &= 6x + 2 \\
    2(x - 3) &= 4 & 3(x + 4) &= -6 & 5(x + 1) &= 2x - 4 \\
    0.2x &= 4 & \frac{x}{3} - 2 &= -4 & 0.5x + 3 &= 1 
\end{align*}
\section{式の展開(式を展開せよ)}
\begin{align*}
    5(x+3) &= & (x+2)(x+6) &= \\
    -2(4x-7) &= & (x-3)(x+5) &= \\
    (x+8)^2 &= & (x-9)^2 &= \\
    (x+12)(x-12) &= & (2x+3)(x+2) &= \\
    (3x-1)(2x+5) &= & (4x-3)^2 &= 
\end{align*}
\section{因数分解(因数分解せよ)}
\begin{align*}
    2x+10 &= & x^2+7x+10 &= \\
    x^2-2x-8 &= & x^2-5x+4 &= \\
    x^2+6x+9 &= & x^2-25 &= \\
    x^2+x-20 &= & x^2+10x+25 &= \\
    x^2-81 &= & x^2+13x+42 &= 
\end{align*}
\section{文字と式(文章を元に式を作成せよ)}
\begin{align*}
　&1.　a円の商品を5個、b円の商品を3個買った。合計金額はc円である。\\
　&2.　1000円札を出して、x円の商品を3個買った。おつりはy円である。\\
　&3.　分速v\,\text{m}でt分間歩いた。進んだ道のりはd\,\text{m}である。\\
　&4.　100\,\text{L}入った水槽から、毎分x\,\text{L}ずつt分間水を抜いた。残りの水の量はy\,\text{L}である。\\
　&5.　x円の商品2個とy円の商品5個の合計金額に、10\%の消費税を加えるとz円になる。\\
　&6.　4回のテストの点数がそれぞれa, b, c, d点だった。平均点はm点である。\\
　&7.　現在地xから、負の方向に速さvでt秒間移動した。到達した地点の座標はpである。\\
　&8.　定価a円の商品をx\%引きで10個買った。合計金額はb円である。
\end{align*}
\end{document}


